\section{Prompt size, token use, and cost}
\label{sec:economics}

A smaller prompt does not necessarily make successful execution cheaper.
For model-written state without retained history, the median peak-prompt
ratio is \ModelStatePeakRatio{}, but the cost ratio is
\ModelStateCostRatio{} on shared successes
(Table~\ref{tab:resources}). The corresponding output-token ratio is
\ModelStateOutputRatio{}. The compact representation therefore coexists
with higher output use and higher reference-tariff cost in this subset.

\begin{table}[htbp]
\centering\small
\caption{Resource ratios on shared successes: median of matched
condition/full-history ratios. Peak prompt is a per-call maximum; the
token columns are cumulative. The exact sign test concerns cost ratios.}
\label{tab:resources}
\input{tables/new-resources}
\end{table}

Retaining history brings the conditional cost ratio toward one:
\ModelHistoryTwoCostRatio{} at $k=2$ and
\ModelHistorySixCostRatio{} at $k=6$. These ratios do not describe the
whole workload. They omit both failures and successes that only one
condition achieves, precisely the outcomes on which capability differs.

Table~\ref{tab:economics} includes retained failures and reports cost and
successes per dollar. Taking the reciprocal of the latter gives cost per
observed success. State without history costs
\$\ModelStateCostPerSuccess{} per success, compared with
\$\ModelControlCostPerSuccess{} for full history. At $k=6$, the value is
\$\ModelHistorySixCostPerSuccess{}, below that reference despite the
higher conditional cost ratio. There is no contradiction: the conditional
comparison measures resources on shared successes, while the all-execution
comparison includes changes in success and termination.

\begin{table}[htbp]
\centering\small
\caption{Cost and efficiency over retained executions, including failures.
Successes per dollar use aggregate cost and observed successes, not ratios
of medians. Token efficiency is reported where
cumulative totals are available. Cost conventions differ between panels.}
\label{tab:economics}
\input{tables/new-economics}
\end{table}

Figure~\ref{fig:pareto} places these history-retention policies in the
success--cost plane. The panels keep their own controls and accounting
conventions, so horizontal positions are compared only within a panel.
% Adapted from the repository's original trade-off scatter.
\begin{figure}[tbp]
\centering
\begin{tikzpicture}
\begin{axis}[
width=.47\linewidth,height=5.1cm,
title={Model-written memory},title style={font=\small},
xlabel={Mean cost (\$)},ylabel={Solved rate (\%)},
xmin=.040,xmax=.078,ymin=30,ymax=68,
xtick={.04,.05,.06,.07},scaled x ticks=false,
xticklabel style={/pgf/number format/fixed,/pgf/number format/precision=2,font=\scriptsize},
yticklabel style={font=\scriptsize},label style={font=\small},
grid=major,grid style={black!12},
nodes near coords={\pgfplotspointmeta},point meta=explicit symbolic,
every node near coord/.append style={font=\scriptsize,anchor=south,yshift=3pt}]
\addplot[only marks,mark=*,mark size=2.5pt,blue!65!black]
table[col sep=comma,x=mean_cost_usd,y=solved_pct,meta=history]{../paper/generated/new-model-history.csv};
\end{axis}
\begin{axis}[
at={(.52\linewidth,0)},anchor=south west,
width=.47\linewidth,height=5.1cm,
title={Hybrid memory},title style={font=\small},
xlabel={Mean cost (\$)},
xmin=.016,xmax=.028,ymin=30,ymax=68,
xtick={.016,.020,.024,.028},scaled x ticks=false,
xticklabel style={/pgf/number format/fixed,/pgf/number format/precision=3,font=\scriptsize},
yticklabel style={font=\scriptsize},label style={font=\small},
grid=major,grid style={black!12}]
\addplot[only marks,mark=square*,mark size=2.5pt,green!45!black]
table[col sep=comma,x=mean_cost_usd,y=solved_pct,meta=history]{../paper/generated/new-hybrid-history.csv};
\node[font=\scriptsize,green!45!black,anchor=south west,xshift=3pt,yshift=3pt]
 at (axis cs:\HybridControlCostMean,\HybridControlRate) {full};
\node[font=\scriptsize,green!45!black,anchor=north east,xshift=-3pt,yshift=-3pt]
 at (axis cs:\HybridHistorySixCostMean,\HybridHistorySixRate) {6};
\node[font=\scriptsize,green!45!black,anchor=north east,xshift=-3pt,yshift=-3pt]
 at (axis cs:\HybridCostMean,\HybridRate) {0};
\node[font=\scriptsize,green!45!black,anchor=south west,xshift=3pt,yshift=3pt]
 at (axis cs:\HybridHistoryTwoCostMean,\HybridHistoryTwoRate) {2};
\end{axis}
\end{tikzpicture}
\caption{Success--cost trade-offs when retaining history. Labels are the retention setting $k$; full denotes the full-history reference without a state block. Each panel includes failures in mean cost and uses its own control and cost convention. The horizontal scales differ.}
\label{fig:pareto}
\end{figure}


Across retained executions, model-written state alone has a lower median
peak prompt but higher median output and more recorded turns than full
history (Table~\ref{tab:tokens}). The largest history-retention setting
increases uncached input, while bringing output closer to the full-history
level. Cache-read counts differ substantially across policies. These
observations establish resource differences, but do not isolate cache
invalidation, repeated inspection, or patch generation as their causes.

\begin{table}[htbp]
\centering\small
\caption{Median recorded turns $T$, per-execution peak prompt, and
cumulative token use for model-written state policies. Token counts are in
thousands; input excludes cache reads. Peak is the median of execution
maxima, not the median per-call prompt. Token columns are separate marginal
medians and cannot be summed to obtain median total tokens.}
\label{tab:tokens}
\input{tables/new-tokens}
\end{table}

The economic objective determines which comparison matters. Peak prompt
size concerns context capacity; cost per success concerns the retained
workload; conditional ratios compare resources when both policies finish.
An evaluation based only on compression can miss a capability deficit,
while one based only on aggregate expense can favor cheap unsuccessful
executions. Reporting these quantities together makes the trade-off
visible without treating any single measure as a universal efficiency
score.
